% Copyright (c) 2026 Christian Reitwiessner.
% Authors: Christian Reitwiessner
% Build: /tmp/andreev-pdf-tools/tectonic --keep-logs paper.tex
\documentclass[11pt]{article}
\usepackage[T1]{fontenc}
\usepackage{lmodern,amsmath,amssymb,amsthm,booktabs,microtype,enumitem}
\usepackage[margin=1in]{geometry}
\usepackage[colorlinks=true,allcolors=blue]{hyperref}
\newtheorem{theorem}{Theorem}[section]
\newtheorem{lemma}[theorem]{Lemma}
\newtheorem{proposition}[theorem]{Proposition}
\newtheorem{corollary}[theorem]{Corollary}
\newcommand{\bits}{\{0,1\}}
\newcommand{\ind}[1]{[\,#1\,]}
\newcommand{\THR}{\mathsf{THR}}
\newcommand{\smalltitle}[1]{\paragraph{#1}}
\setlength{\emergencystretch}{1em}
\setlist[itemize]{itemsep=1pt,topsep=3pt}
\setlist[enumerate]{itemsep=1pt,topsep=3pt}
\title{Quadratic Depth-Two Threshold Circuits\\for Andreev's Function}
\author{Christian Reitwiessner}
\date{November 2026}
\begin{document}
\maketitle
\begin{abstract}
Andreev's function reads an arbitrary table at an address obtained by taking
parities of disjoint input blocks. We give an exact depth-two threshold
circuit with $(5/16+o(1))n^2$ gates, improving the previously stated
$O(n^3/\log n)$ upper bound. For uniform inputs and $0<\delta\le1/2$, error $\delta$ needs
$O(n^{3/2}\sqrt{\log(1/\delta)})$ gates; a tunable version attains advantage
$\varepsilon$ with $O(\varepsilon^2n^{3/2}\sqrt{\log(1/\varepsilon)})$ gates
for $N^{-1/2}\le\varepsilon\le1/8$, where $N$ is the table length.
For fixed $0<\delta<1/2$, the upper exponent matches known lower bounds up to
logarithms. We also prove an $\Omega(n^2/\log^3n)$ lower bound when every
bottom gate has weights of one sign within each address block. The exact
construction belongs to this class. In a different restricted class, where
each nonzero address normal has full support and equal absolute weights,
we prove an exponential lower bound. The proofs use Fourier inversion,
threshold staircases, a cube-counting argument, and literal identifications.
\end{abstract}

\section{Introduction}
A depth-two threshold circuit consists of threshold gates on the input bits
and one threshold gate combining their outputs. The output gate may also
read the raw inputs. All real weights are allowed. This simple model still
has a substantial gap between known upper and lower bounds for explicit
functions. We study the ordinary lookup version of Andreev's function.

Fix $k,m\ge1$ and let $N=2^k$. The input is an $N$-bit table
$x=(x_c)_{c\in\bits^k}$ and $k$ blocks $a_i\in\bits^m$. Define
$$b_i=|a_i|\bmod2,\qquad A_{k,m}(x,a)=x_b,\qquad n=N+km.$$
Our asymptotic statements use $m=\lfloor N/k\rfloor$, so $km\sim N$ and
$n\sim2N$. When exact balance $km=N$ is used explicitly, integrality
requires $k\mid2^k$; the floor-balanced family avoids this restriction.
The table entries remain independent input bits throughout.

\paragraph{Results and proof ideas.}
The exact circuit has $(5/16+o(1))n^2$ gates and polynomial integer weights
(Theorem~\ref{thm:exact}). Fourier inversion writes lookup as a sum of $N$
products of a table linear form and an address parity. We implement each
whole product with staircases on raw affine forms. A sign-only refinement
shares the quadratic correction terms and handles a few exceptional
addresses with a margin argument. No intermediate parity is computed in
an extra circuit layer.

For uniform independent table and address bits, Theorem~\ref{thm:avgupper}
gives an $O(n^{3/2})$-gate circuit at every fixed error. Independent fixed
address masks cancel cross-character truncation errors in mean square.
Using only the characters on a smaller set of blocks gives the stated
size--advantage tradeoff. The masks are fixed nonuniformly by averaging;
they are not chosen separately for each input.

A gate is \emph{sign-coherent} if its address weights in each block are all
nonnegative or all nonpositive; orientations may differ between blocks
and gates. Theorem~\ref{thm:coherent} gives $\Omega(n^2/\log^3n)$ bottom
gates for this class, with arbitrary table and top weights. Thus our
quadratic circuit is optimal within this class up to logarithms. The proof
runs through every unit cube of the block-count grid: each cube reveals
the entire table, while a coherent gate changes on few cubes.
Theorem~\ref{thm:fullsign} gives a separate exponential obstruction for
full-support equal-magnitude address normals. Zero supports and unequal
magnitudes are essential qualifications, not harmless technical details.

\paragraph{Related work and attribution.}
Kane and Williams~\cite{KW} proved $\widetilde\Omega(n^{3/2})$ gate and
$\widetilde\Omega(n^{5/2})$ wire lower bounds for this function, including
average-case bounds. Their Theorem~1.2 and Section~5.2 give the
$O(n^3/\log n)$ depth-two upper bound, and a quadratic construction in the
different $\mathsf{MOD}_3\circ\mathsf{MOD}_2$ basis. They Fourier-expand the
lookup tests and separately simulate the resulting parities. Fourier
inversion itself is standard; the savings here come from grouping terms
by address character before compiling them. The staircase representation
of a bounded integer linear form is also standard; see the symmetric-function
construction of Hajnal et al.~\cite[Proposition~2.1]{HMPST} and the related
representation theorem of Hansen and Podolskii~\cite{HP}.
Kane--Williams use concentration-window staircases for approximate parity;
our approximation proof combines that tool with lookup and cancellation
of cross-character errors.

Impagliazzo, Paturi and Saks~\cite{IPS} use restrictions and monotone-chain
arguments for threshold lower bounds. Our coherent proof combines the
single-crossing property with an information count over all count-grid
cubes. Kabanets, Kane and Lu~\cite{KKL} extend restriction methods to
polynomial threshold functions. Chen, Santhanam and Srinivasan~\cite{CSS}
study average-case threshold hardness, including a \emph{generalized}
Andreev function with encoded rather than independent table entries.
These are distinct parameterized functions, as is polynomial-containment
``Andreev's Problem''~\cite{Pot}.
We make no exhaustive priority claim for the elementary identities or
restricted-class corollaries in this note.

\section{Two identities for implementing lookup}\label{sec:identities}
We count bottom gates; one output gate is added to every upper bound.
The notation $[L\ge t]$ denotes the Boolean value of the comparison, so
any such comparison is a single bottom gate when $L$ is affine in raw
inputs. An affine score is free at the top. Wires count nonzero aggregated
edges, not coefficient multiplicities. All logarithms are natural except
$\log_2$ and information entropies.

In Sections~\ref{sec:identities}--\ref{sec:approx}, use
$$X_S=\sum_{c\in\bits^k}(-1)^{S\cdot c+x_c},\qquad
 t_S=\sum_{i\in S}|a_i|,\qquad S\cdot c=\sum_{i\in S}c_i.$$
Here $S\subseteq[k]$, $X_S$ is even and lies in $[-N,N]$, and
$0\le t_S\le m|S|$. The subscript $S$ always indexes a set of blocks.
\begin{lemma}[Lookup as a sum of products]\label{lem:lookup}
$\displaystyle\sum_{S\subseteq\{1,\dots,k\}}X_S\,(-1)^{t_S}=N\,(-1)^{x_b}.$
\end{lemma}
\begin{proof}
Since $|a_i|\equiv b_i\pmod2$, we have $(-1)^{t_S}=(-1)^{S\cdot b}$.
Exchanging the two sums, the left side is
$$
\sum_c(-1)^{x_c}\prod_{i=1}^k\bigl(1+(-1)^{c_i+b_i}\bigr),
$$
and the product is $N$ for $c=b$ and $0$ otherwise.
\end{proof}

This identity is also a change of basis. The truth-table vector
$v_c=(-1)^{x_c}$ has unnormalized Fourier coordinates $X=Hv$, where
$H_{S,c}=(-1)^{S\cdot c}$ is the Walsh--Hadamard matrix. Its columns
are orthogonal: $H^{\mathsf T}H=NI$. Thus $v=N^{-1}H^{\mathsf T}X$,
and Lemma~\ref{lem:lookup} is precisely its $b$th coordinate, since
$(-1)^{t_S}=H_{S,b}$. Equivalently, $X_S/N$ is the coefficient of the
parity function $c\mapsto(-1)^{S\cdot c}$ in the Fourier expansion of
$c\mapsto(-1)^{x_c}$.

Let $\varphi(u)=(-1)^u u$ for integers $u$.

\begin{lemma}[Removing the parity factor]\label{lem:parity}
If $X$ is an even integer and $t$ is an integer, then
$X(-1)^t=\varphi(X+t)-\varphi(t)$.
\end{lemma}
\begin{proof}
$X+t$ has the parity of $t$, so $\varphi(X+t)=(-1)^t(X+t)$.
\end{proof}

\begin{lemma}[Threshold staircase]\label{lem:ladder}
Let $L$ be an integer that always lies in $\{\ell,\dots,h\}$, and let $g$
be any function on the integers. Then
$$
g(L)=g(\ell)+\sum_{\theta=\ell+1}^{h}\bigl(g(\theta)-g(\theta-1)\bigr)\ind{L\ge\theta}.
$$
\end{lemma}
\begin{proof}
The brackets equal $1$ exactly for $\theta\le L$, and the sum telescopes.
\end{proof}

When $L$ is a linear form in the input bits, each $\ind{L\ge\theta}$ is one
threshold gate. Both $t_S$ and $X_S+t_S$ are such forms, since
$(-1)^{x_c}=1-2x_c$:
\begin{equation}\label{eq:form}
X_S+t_S=\sum_c(-1)^{S\cdot c}-2\sum_c(-1)^{S\cdot c}x_c+\sum_{i\in S}\sum_{j=1}^m a_{i,j}.
\end{equation}


For example, apply the last two lemmas to
$$N(-1)^{x_b}=\sum_S\bigl(\varphi(X_S+t_S)-\varphi(t_S)\bigr).$$
The two ranges have lengths $2N+m|S|$ and $m|S|$. Since
$\sum_S|S|=kN/2$, this already gives $2N^2+kmN$ bottoms, or
$3n^2/4$ at exact balance. The next section reduces the leading
constant by asking only for a correct sign.

\section{The sharper exact construction}\label{sec:exact}
The proof has two independent checks: first the sign of an integer score,
then the raw threshold gates realizing it. Its auxiliary notation is local
to this section; none of it is used in the lower bounds.
\begin{theorem}\label{thm:exact}
Put $\alpha=\lfloor(m-1)/2\rfloor$. There is an exact depth-two threshold
circuit for $A$ with at most
\begin{equation}\label{eq:bestcount}
\frac{N(N-1)}2+\frac{kN(m+\alpha)}2
 +2(N-1)+2\alpha k+3km+\left\lceil\frac{km}{2}\right\rceil
\end{equation}
bottom gates and one top gate. Its weights are integers of magnitude
polynomial in $n$. At exact balance its total size is
$$
\frac5{16}n^2-\frac n4\log_2(n/2)+O(n).
$$
For $m=\lfloor N/k\rfloor$, its size is $(5/16+o(1))n^2$.
\end{theorem}

\begin{proof}
\emph{Forms used only in this proof.}
Write $s_i=|a_i|$, $\sigma=\sum_i s_i$, and let $e_i$ be the address with
only coordinate $i$ equal to one; $x_0$ means the all-zero address entry. Put
$$
z=\sum_{c\ne0}x_c,\quad E=\sum_i x_{e_i},\quad
w=2z+2\alpha E+\sigma,\quad R=2(N-1)+2\alpha k+km,\quad C=N(R+4).
$$
Thus $0\le w\le R$. For each nonempty $S\subseteq[k]$, define
$$
W_S=\sum_{c:S\cdot c\text{ odd}}x_c+\alpha\sum_{i\in S}x_{e_i},\quad
\kappa_S=-4(2\alpha+1)(k-2|S|),
$$
and let $r_S(u)=-8\lceil u/2\rceil^2+\kappa_S\lceil u/2\rceil$.
All of $W_S,t_S,w,s_i$ are affine forms in the original inputs.

\emph{The score.} Consider
\begin{align}\label{eq:bestscore}
P={}&2N+\sum_{S\ne\varnothing}r_S(2W_S+t_S)
 +\frac N2\left(w^2+\sum_i s_i^2\right)\notag\\
 &+4Nz-2N(\alpha^2+1)E-2\alpha N\sigma
 +2N(1+\alpha)\sum_i x_{e_i}s_i-Cx_0\ind{b=0}.
\end{align}
We first verify its sign, then realize its non-affine terms with bottom gates.
The following identities follow by expanding the squares and summing
characters; all sums over $S$ are over nonempty sets:
\begin{align*}
\sum_S4W_S(-1)^{t_S}
 &=2N\ind{b=0}z-2N\ind{b\ne0}x_b
   +2\alpha N\sum_i x_{e_i}\bigl(\ind{b=0}-\ind{b=e_i}\bigr),\\
8\sum_S(W_S^2+W_St_S)
 &=\frac N2(w^2-\sigma^2)+2Nz+(4\alpha+2\alpha^2)NE
   +2N(1+\alpha)\sum_i x_{e_i}s_i,\\
8\sum_S\lceil t_S/2\rceil^2
 &=\frac N2\left(\sigma^2+\sum_i s_i^2\right)
   +N(\sigma+1)\ind{b\ne0}+N\sum_i s_i\ind{b=e_i}.
\end{align*}
To check the quadratic expansion, for distinct nonzero addresses $c,d$ use
$$\sum_S[S\cdot c\text{ odd}]=N/2,\qquad
\sum_S[S\cdot c\text{ odd}][S\cdot d\text{ odd}]=N/4.$$
For $c=d$ the second sum is $N/2$. These counts expand $W_S^2$ directly;
the other terms use the same orthogonality and
$\lceil t/2\rceil=(t+[t\text{ odd}])/2$.
Also
$$
4W(-1)^t=4W-8\lceil(2W+t)/2\rceil^2+8W^2+8Wt+8\lceil t/2\rceil^2.
$$
The shifts $\kappa_S$ change only coefficients of existing odd thresholds,
so the singleton-address correction adds no staircase gates. They supply
the address correction because
$$
\sum_S\kappa_S\bigl(\lceil t_S/2\rceil-t_S/2\bigr)
 =N(2\alpha+1)\sum_i\ind{b=e_i}.
$$
The remaining linear shifts satisfy
$$\sum_S\kappa_SW_S=2N(2\alpha+1)(1+\alpha)E,\qquad
\frac12\sum_S\kappa_St_S=N(2\alpha+1)\sigma.$$
These follow by writing $k-2|S|=\sum_i(-1)^{[i\in S]}$ and applying the
same orthogonality. Substitution in \eqref{eq:bestscore} gives
\begin{equation}\label{eq:bestclosed}
P=\begin{cases}
N(w+2)-Cx_0,&b=0,\\
2N(1+\alpha)(1-x_{e_i})-Ns_i,&b=e_i,\\
N(1-2x_b),&\text{otherwise}.
\end{cases}
\end{equation}
At $b=0$, use $0\le w\le R$ and $C=N(R+4)$. At $b=e_i$, the odd
integer $s_i$ lies between $1$ and $2\alpha+1$. In every case $P\ge N$
when $x_b=0$, and $P\le-N$ when $x_b=1$. The top gate outputs $\ind{-P\ge0}$.

\emph{Raw gates and their count.}
The staircase for $r_S(2W_S+t_S)$ uses only odd thresholds
$$
\ind{2W_S+t_S\ge\theta},\qquad
1\le\theta\le N+2\alpha|S|+|S|m,
$$
with coefficient $-8\theta+\kappa_S$. The two square terms use $R+km$
gates by Lemma~\ref{lem:ladder}. For $B=n+2$, each product $x_{e_i}s_i$
uses $m$ gates $\ind{s_i+Bx_{e_i}\ge B+\ell}$, $1\le\ell\le m$.
Finally, the identity
$$
x_0\ind{b=0}=x_0-x_0\sigma+
 \frac4N\sum_{S\ne\varnothing}x_0\lfloor t_S/2\rfloor
$$
implements the last correction with $km$ gates
$\ind{\sigma+Bx_0\ge B+\ell}$ and
$\sum_S\lfloor |S|m/2\rfloor$ gates
$\ind{t_S+Bx_0\ge B+2j}$. Its identity follows from
$\sum_S\ind{t_S\text{ odd}}=(N/2)\ind{b\ne0}$.
The even-level $\sigma$ gates coincide with the $S=[k]$ gates, saving
$\lfloor km/2\rfloor$ gates. Everything else in $P$ is an affine top term.

Since $\sum_{S\ne\varnothing}|S|=kN/2$, the main-ridge and last-correction
counts together are
$$
\sum_{S\ne\varnothing}\left(
 \left\lceil\frac{N+2\alpha|S|+|S|m}{2}\right\rceil
 +\left\lfloor\frac{|S|m}{2}\right\rfloor\right)
 =\frac{N(N-1)}2+\frac{\alpha kN}2+\frac{kmN}2.
$$
Adding the remaining gates gives \eqref{eq:bestcount}. Vanishing coefficients
or further coincidences can only improve it. At exact balance $m$ is even,
so $\alpha=m/2-1$; substitution gives the stated asymptotic bound.
For floor balance, the rounding changes the leading expression by
$O(Nk)=o(n^2)$. Bottom coefficients and thresholds are $O(n)$; the displayed
expansion has top coefficients bounded by $O(n^3)$.
\end{proof}


\section{Approximation by masked staircases}\label{sec:approx}
We now return to $X_S,t_S$ from Section~\ref{sec:identities}; the auxiliary
variables of the exact construction are no longer in use. Agreement is
always under the uniform distribution on all input bits. For this section,
``balanced'' means $m=\lfloor N/k\rfloor$ or $km=N$; only $km=O(N)$ and
$n=\Theta(N)$ enter the estimates.
\begin{theorem}\label{thm:avgupper}
In the balanced case, the following depth-two threshold circuits exist:
\begin{enumerate}
\item For $0<\delta\le1/2$, agreement at least $1-\delta$ with
      $O\bigl(n^{3/2}\sqrt{\log(1/\delta)}\bigr)$ gates.
\item For $N^{-1/2}\le\varepsilon\le1/8$, agreement at least
      $1/2+\varepsilon$ with
      $O\bigl(\varepsilon^2n^{3/2}\sqrt{\log(1/\varepsilon)}\bigr)$ gates. Its support-wire count is\\
      $O\bigl(\varepsilon^2n^{5/2}\sqrt{\log(1/\varepsilon)}\bigr)$.
\end{enumerate}
For $1/8<\varepsilon<1/2$, Part~1 applies with $\delta=1/2-\varepsilon$.
The proof chooses one fixed circuit for each size; it does not assert an
efficient uniform algorithm for choosing its masks.
\end{theorem}

\begin{proof}
\emph{A majority vote with fewer characters.}
Fix $0\le r\le k$ and put $q=2^r$. Character orthogonality gives
$$
P_*=2\sum_{S\subseteq[r]}X_S(-1)^{t_S}
 =2q\sum_{c:\,c_{[r]}=b_{[r]}}(-1)^{x_c}.
$$
Its sign is a majority vote of $L=N/q$ independent table bits, one of them
$x_b$. With ties outputting zero, its agreement is
$1/2+p(L)/2$, where $p(1)=1$ and
$$
p(L)=2^{-L}\binom{L}{L/2}\ge(2L)^{-1/2}\qquad(L\ge2).
$$
For even $L$, the binomial expression satisfies
$p(L+2)/p(L)=(L+1)/(L+2)\ge\sqrt{L/(L+2)}$; together with
$p(2)=1/2$, induction proves the inequality.
For $L\ge2$, $L$ is even because it is a power of two. Conditioning on
the selected bit leaves $L-1$ independent signs; the excess agreement is
half the probability that their sum equals $-1$. The case $L=1$ is exact.
The vote is $[-P_*\ge2q]$ if $q<N$, and $[-P_*\ge0]$ if $q=N$.

\emph{Masked, truncated staircases.}
Independently for each nonempty $S\subseteq[r]$, choose a uniform mask $z_S$
on its $M_S=|S|m$ active raw address coordinates. Let
$$
D_S=|a_S\oplus z_S|-\lfloor M_S/2\rfloor,\qquad
\eta_S=(-1)^{|z_S|+\lfloor M_S/2\rfloor}.
$$
With $\varphi(u)=(-1)^u u$, evenness of $X_S$ gives
$$
\eta_S\bigl(\varphi(X_S+D_S)+\varphi(X_S-D_S)\bigr)
 =2X_S(-1)^{t_S}.
$$
For an integer radius $R\ge1$, set
$f_R(u)=\varphi(u)\ind{|u|\le R}$, and replace the left side by its
$f_R$ version. Together with the free affine term $2X_\varnothing$, these
terms give a score $P_z$. The finite-difference expansion
$$
f_R(U)=\sum_{\ell=-R}^{R+1}
 \bigl(f_R(\ell)-f_R(\ell-1)\bigr)\ind{U\ge\ell}
$$
uses at most $2R+2$ gates. Each $X_S\pm D_S$ is an affine form in the raw
inputs, with table weights $\pm2$ and address weights $\pm1$.
Thus $P_z$ uses at most $4(q-1)(R+1)$ bottoms and one top gate.

\emph{Cancellation of cross terms.}
Let $e_S$ be the truncated contribution minus its exact contribution.
Conditional on $x,a$, the masks are independent, and
$D_S$ has a shifted binomial law independent of $a$. Moreover
$\eta_S=(-1)^{t_S}(-1)^{D_S}$. Therefore
$$
\mathbb E_{z_S}[e_S\mid x,a]=(-1)^{t_S}g_S(X_S)
$$
for a function $g_S:\mathbb Z\to\mathbb R$ independent of $a$. For $S\ne T$, averaging the product
over masks and then uniform raw addresses gives
$$
\mathbb E_{z,x,a}[e_Se_T]
 =\mathbb E_x[g_S(X_S)g_T(X_T)]\,
   \mathbb E_a[(-1)^{t_S+t_T}]=0.
$$
No independence of Walsh coefficients is assumed. Consequently
\begin{equation}\label{eq:maskl2}
\mathbb E_{z,x,a}(P_z-P_*)^2=\sum_{S\ne\varnothing}\mathbb E e_S^2.
\end{equation}

\emph{Diagonal terms and one fixed circuit.}
For $U=X_S\pm D_S$, the mean has magnitude at most $1/2$ and the centered
variable is subgaussian with variance proxy $v_S=N+M_S/4$. Integrating its
Hoeffding tail yields
$$
\mathbb E\bigl[U^2\ind{|U|>R}\bigr]
 \le16v_S\exp\bigl(-R^2/(16v_S)\bigr).
$$
Here are the constants explicitly. Put $V=U-\mathbb EU$. Since
$|\mathbb EU|\le1/2$ and $R\ge1$, $|U|>R$ implies $|V|>R/2$ and
$U^2\le2V^2+1/2$. Integrating the tail
$\Pr(|V|>t)\le2e^{-t^2/(2v_S)}$ gives
$$\mathbb E[U^2[|U|>R]]\le(R^2+8v_S+1)e^{-R^2/(8v_S)}
 \le16v_Se^{-R^2/(16v_S)}.$$
The last step uses $v_S\ge2$ and $1+t\le e^t$. Put $v=N+rm/4$. Since zero truncation removes only tail values,
\eqref{eq:maskl2} implies
$$
\mathbb E_{z,x,a}(P_z-P_*)^2
 \le64(q-1)v\exp\bigl(-R^2/(16v)\bigr).
$$
If $q=N$, the untruncated score is $\pm2N$, so test $-P_z\ge0$.
If $q<N$, its attainable values are multiples of $4q$; testing
$-P_z\ge2q$ reproduces its negative-sign vote, with a gap of $2q$ on
either side. In both cases disagreement with the vote requires
$|P_z-P_*|\ge2q$. Markov's inequality gives expected disagreement at most
$$
\frac{16v}{q}\exp\bigl(-R^2/(16v)\bigr).
$$
Taking
$$
R=\left\lceil4\sqrt{v\log\frac{16v}{q\delta}}\right\rceil
$$
makes this at most $\delta$. Averaging fixes one mask collection with that
uniform-input error. The resulting circuit is deterministic; masks are
neither supplied inputs nor chosen separately for different inputs.

\emph{The two parameter choices.}
For Part~1 take $q=N$. At balance $v=O(N)$, so
$R=O(\sqrt{N\log(1/\delta)})$. For Part~2 take the least power of two
$q\ge32\varepsilon^2N$ and set $\delta=\varepsilon$.
Then $q<64\varepsilon^2N$, $q\le N$, and the vote's advantage is at least
$\sqrt{q/(8N)}\ge2\varepsilon$. Its truncated circuit has the required
advantage, with $\log(16v/(q\delta))=O(\log(1/\varepsilon))$.
With $G$ bottoms the full wire count is at most $nG+G+n$, including
raw top wires and gate-to-top edges. The parameter range absorbs the
additive terms, proving the stated wire bound.
\end{proof}
\section{A near-quadratic bound for coherent gates}\label{sec:coherent}
This section uses only the original parameters $N,k,m$, the table $x$, and
a bottom-gate count $G$. A coherent gate may have arbitrary real table
weights, and the top gate is unrestricted. For the next two lemmas put
$$D_0=\frac{\max\{0,N-k-1-\log_2(G+1)\}}{k+\log_2(G+1)}.$$
\begin{theorem}\label{thm:coherent}
Every exact circuit whose bottoms are sign-coherent satisfies
$$G\ge\frac{D_0}{1-(1-1/m)^k}\ge\frac m k D_0.$$
In particular, for $m=\lfloor N/k\rfloor$ and sufficiently large $k$,
$G=\Omega(n^2/\log^3n)$. The same holds at exact balance.
\end{theorem}
The upper construction in Section~\ref{sec:exact} belongs to this class.
We prove the bound by comparing how many gates must change on a cube
with how many cubes one coherent gate can cross.
\smalltitle{The cube restriction.}
Fix an order of the positions inside every block. For a vector
$s\in\{0,\dots,m\}^k$ let $a(s)$ be the address part in which block $i$ has
ones exactly at its first $s_i$ positions. A \emph{cube} is a vector
$q\in\{0,\dots,m-1\}^k$; there are $m^k$ cubes. For $z\in\bits^k$ the input
$(x,a(q+z))$ differs from $(x,a(q))$ exactly in the live positions
$q_i+1$ of the blocks $i$ with $z_i=1$, and its address is
$(q\bmod2)\oplus z$. Hence
$$
z\mapsto A\bigl(x,a(q+z)\bigr)=x_{(q\bmod 2)\oplus z}
$$
determines the whole table $x$. A bottom gate \emph{splits} the cube $q$ for
the table $x$ if it is non-constant in $z$ on the inputs $(x,a(q+z))$. Let
$d_q(x)$ be the number of bottom gates that split $q$ for $x$.

\begin{lemma}[Many gates split each cube]\label{lem:entropy}
Let $C$ be any depth-two threshold circuit for $A$ with $G$ bottom gates, and
let $x$ be uniformly random. Then $\mathbb E_x[d_q(x)]\ge D_0$ for every cube
$q$.
\end{lemma}
\begin{proof}
Fix $q$. On the inputs $(x,a(q+z))$, bottom gate $j$ computes
$\ind{\lambda_j\cdot z\ge\tau_j}$, where $\lambda_j\in\mathbb R^k$ lists its
weights on the live positions and only the number $\tau_j$ depends on $x$.
Likewise the top gate computes
$\ind{\sum_jv_jg_j+\mu\cdot z\ge\kappa}$, where $v_j$ are its weights on the
bottom gates and $\mu$ its weights on the live positions; the gates that
do not split $q$ are constant and are absorbed, together with the table and
the fixed address bits, into $\kappa$. After absorption the sum is over the splitting set $I$ only.
Thus only $\tau_j$ and $\kappa$ depend on $x$.

Therefore the function $z\mapsto x_{(q\bmod2)\oplus z}$, and with it $x$, is
determined by: the set $I$ of splitting gates; for each $j\in I$ the
non-constant function $\ind{\lambda_j\cdot z\ge\tau_j}$, one of at most $2^k-1$
(only the position of $\tau_j$ among the $\le2^k$ values $\lambda_j\cdot z$
matters); and the position of $\kappa$, one of at most $2^k+1$. With
$D=|I|=d_q(x)$, entropy gives
$$
N=H(x)\le H(D)+\mathbb E\log_2\Bigl[\tbinom GD(2^k-1)^D(2^k+1)\Bigr]
\le\log_2(G+1)+\mathbb E[D]\,(\log_2(G+1)+k)+k+1,
$$
which rearranges to the claim; a negative numerator is replaced by the
trivial lower bound zero.
\end{proof}

\begin{lemma}[A coherent gate splits few cubes]\label{lem:diag}
For every table $x$, a sign-coherent gate splits at most $m^k-(m-1)^k$ cubes.
\end{lemma}
\begin{proof}
Let $\varepsilon_i=+1$ if the gate's weights in block $i$ are $\ge0$ and
$\varepsilon_i=-1$ otherwise. Its linear form on $(x,a(s))$ is a constant plus
$\sum_iF_i(s_i)$, where $F_i(s_i)$ is the sum of its weights on the first
$s_i$ positions of block $i$. So the form $F(s)$ does not decrease when
$s_i$ moves by one step in direction $\varepsilon_i$.

On the cube $q$ the form is smallest at $q^-=q+(\mathbf1-\varepsilon)/2$ and
largest at $q^+=q+(\mathbf1+\varepsilon)/2$. The gate with threshold $\theta$
splits $q$ exactly when $F(q^-)<\theta\le F(q^+)$.

Group the cubes into \emph{lines} $q,\,q+\varepsilon,\,q+2\varepsilon,\dots$.
For consecutive cubes on a line, the largest vertex of one is the smallest
vertex of the next, so the intervals $(F(q^-),F(q^+)]$ along a line are
disjoint, and $\theta$ lies in at most one of them. Each line starts at a cube
$q$ for which $q-\varepsilon$ is not a cube, and exactly $(m-1)^k$ cubes $q$
have $q-\varepsilon$ a cube. So there are $m^k-(m-1)^k$ lines.
\end{proof}

\begin{proof}[Proof of Theorem~\ref{thm:coherent}]
Sum Lemma~\ref{lem:entropy} over all cubes and compare with
Lemma~\ref{lem:diag}: $m^kD_0\le\sum_q\mathbb E[d_q]\le G\,(m^k-(m-1)^k)$.
Then use $1-(1-1/m)^k\le k/m$. If $G\ge N^2$ there is nothing to prove. Otherwise $\log_2(G+1)\le2k$ and
$D_0\ge(N-3k-1)/(3k)$. For $m=\lfloor N/k\rfloor$, we have $m=\Theta(N/k)$ and $n=\Theta(N)$,
giving the claimed asymptotic bound. Exact balance gives the same estimate.
\end{proof}


\section{Full-support equal-magnitude address weights}\label{sec:fullsign}
A different restriction permits mixed signs but requires every nonzero
address coefficient in a gate to have the same absolute value, with no
address coordinates omitted. This is much more restrictive than allowing
arbitrary integer weights. The argument below combines literal identifications with the classical
parity-order obstruction of Minsky and Papert~\cite[Section~3.1]{MP}:
a threshold of predicates computing parity must contain a predicate
depending on every variable. We give the short orthogonality proof inside
the argument, including the treatment of raw top inputs.
\begin{theorem}\label{thm:fullsign}
Let $k\ge2$, $m\ge1$, and $d=m\lfloor k/2\rfloor\ge2$. Suppose a depth-two
threshold circuit computes $A_{k,m}$ exactly, and every bottom gate's address
normal is either zero or of the form
$$
\gamma_j=\rho_j\sigma_j,\qquad \rho_j>0,\quad
\sigma_j\in\{-1,1\}^{km}.
$$
Table coefficients, thresholds, top coefficients, and direct raw-input
coefficients at the top may be arbitrary real numbers. Then the circuit
has at least $2^d$ bottom gates. If $km=\Theta(2^k)$, this is $2^{\Omega(n)}$.
\end{theorem}

\begin{proof}
Put $r=\lfloor k/2\rfloor$. Fix the table to
$x_c=\bigoplus_{i=1}^r c_i$. If $k$ is odd, also fix one block outside
$\{1,\ldots,r\}$ to zero. There remain $2r$ live blocks: the output is
parity on the $d=rm$ bits in the first $r$ blocks and ignores the other
$d$ bits. Write each sign coordinate as $2a-1$, and call the resulting selected
and ignored coordinates $z_i,w_i\in\{-1,1\}$. Pair these selected
and ignored bits arbitrarily. After possibly reversing its orientation,
the circuit's top score has the sign of
$$
f(z,w)=\prod_{i=1}^d z_i.
$$
We may take this score strictly sign-correct on the finite domain: for
output convention $[F\ge0]$, add a sufficiently small positive constant
to $F$ before reversing orientation if necessary. Negative-class scores
were already strictly negative. This preserves all classifications and
all gate counts.

For each $t\in\{-1,1\}^d$, substitute
$$
w_i=t_i z_i\qquad(1\le i\le d).
$$
In bit coordinates these are identifications or complemented
identifications, so every bottom remains a threshold gate and each raw
top term remains affine. The target stays $f(z)=\prod_i z_i$.

For a gate with nonzero address normal, let $\sigma_i,\tau_i$ denote its
signs on the paired selected and ignored coordinates. Its coefficient
on $z_i$ after substitution is a common positive scale times
$\sigma_i+t_i\tau_i$. This is nonzero exactly when
$t_i=\sigma_i\tau_i$. Hence that gate can depend on all $d$ live variables
on at most one of the $2^d$ choices of $t$. A pure-table gate is constant
after fixing the table and serves none of them.

For each fixed $t$, some bottom must depend on every live variable.
Otherwise every bottom misses at least one coordinate, and therefore
has zero correlation with $f$ under the uniform live cube: summing over
that missing coordinate cancels the two signs of $f$. Every affine raw
top term also has zero correlation with $f$, since $d\ge2$. Thus the
whole top score $F_t$ would satisfy $\mathbb E[fF_t]=0$. Strict sign
correctness instead gives
$$
\mathbb E[f(z)F_t(z)]=\mathbb E[|F_t(z)|]>0,
$$
a contradiction. Each of the $2^d$ restrictions therefore needs a bottom
that depends on every live variable, and each bottom can serve at most
one restriction. This proves the bound.
\end{proof}


The full-support assumption cannot be dropped: Section~\ref{sec:identities}
already gives $O(n^2)$ gates for $A$ with address coefficients in $\{0,1\}$.
Exact equality of magnitudes is also essential. Using $X_S,t_S,\varphi$
from Section~\ref{sec:identities}, a full-support quadratic
construction results by giving a character's selected raw coordinates odd
weights of magnitude one and its other coordinates even weights of
magnitude two. The resulting integer form $L_S$ obeys
$(-1)^{L_S}=(-1)^{t_S}$, and the identity
$X_S(-1)^{t_S}=\varphi(X_S+L_S)-\varphi(L_S)$ still applies. Its ranges
are $O(n)$, giving $O(n^2)$ raw gates. Thus the exponential conclusion
asserts no unrestricted lower-bound improvement.

\section{Open questions}
\begin{enumerate}
\item Can exact ordinary Andreev be computed with $n^{2-\eta}$ gates for
some fixed $\eta>0$, or can its unrestricted $\widetilde\Omega(n^{3/2})$
lower bound be improved in exponent? The coherent result requires any
such upper improvement to use mixed signs inside blocks. The exponential
result also rules out relying exclusively on full-support equal-magnitude
address normals.
\item Can the exact wire upper bound be made $o(n^3)$, ideally
$O(n^{3-\eta})$? The known unrestricted wire lower exponent is $5/2$.
The gate upper bound gives $O(n^3)$ wires. Transferring a gate
\emph{lower} bound to wires requires additional support information.
\item What is the optimal size as error tends to zero? The upper bound
$O(n^{3/2}\sqrt{\log(1/\delta)})$ interpolates to quadratic size at
exponentially small error, but the known lower bounds do not give a
matching dependence on $\delta$.
\item Can one efficiently choose the fixed masks in
Theorem~\ref{thm:avgupper}, preserving its bounds? The present existence
proof averages over independently chosen masks.
\end{enumerate}

\paragraph{Proof status and reproducibility.}
The claims here are mathematical proofs. Exact finite checking scripts
accompany the source to check the refined score, raw-gate counts and threshold ties
(\texttt{verify\_results.py}), and masked restrictions, coherent cubes and
pairing identities (\texttt{restriction\_checks.py}). These computations
do not replace the all-size arguments. The earlier $3n^2/4$ construction
has a merged Lean formalization\footnote{\url{https://github.com/cslib-community/cslib-frontier/pull/2813}.},
but this note does not claim formal verification of
the refined constant, the mean-square improvement, or the restricted lower
bounds. The research and exposition were developed with AI assistance.

\begin{thebibliography}{9}\small
\bibitem{KW} D.~M. Kane and R.~Williams.
Super-Linear Gate and Super-Quadratic Wire Lower Bounds for Depth-Two and
Depth-Three Threshold Circuits. STOC, 2016.
\href{https://arxiv.org/abs/1511.07860}{arXiv:1511.07860}.
\bibitem{HMPST} A.~Hajnal, W.~Maass, P.~Pudl\'ak, M.~Szegedy, and G.~Tur\'an.
Threshold circuits of bounded depth. J. Comput. Syst. Sci. 46(2):129--154, 1993.
\bibitem{HP} K.~A. Hansen and V.~V. Podolskii.
Polynomial threshold functions and Boolean threshold circuits.
Information and Computation 240:56--73, 2015.
\bibitem{MP} M.~L. Minsky and S.~A. Papert.
\emph{Perceptrons: An Introduction to Computational Geometry}.
MIT Press, 1969; expanded edition, 1988. Section~3.1.
\bibitem{IPS} R.~Impagliazzo, R.~Paturi, and M.~Saks.
Size--depth tradeoffs for threshold circuits. SIAM J. Comput. 26(3):693--707, 1997.
\bibitem{KKL} V.~Kabanets, D.~M. Kane, and Z.~Lu.
A Polynomial Restriction Lemma with Applications. STOC, 2017.
\href{https://eccc.weizmann.ac.il/report/2017/026/}{ECCC TR17-026}.
\bibitem{CSS} R.~Chen, R.~Santhanam, and S.~Srinivasan.
Average-Case Lower Bounds and Satisfiability Algorithms for Small Threshold
Circuits. Theory of Computing 14(9):1--55, 2018.
\bibitem{Pot} A.~Potukuchi.
On the $\mathsf{AC}^0[\oplus]$ Complexity of Andreev's Problem.
FSTTCS, 2019, Article 25.
\href{https://doi.org/10.4230/LIPIcs.FSTTCS.2019.25}{doi:10.4230/LIPIcs.FSTTCS.2019.25}.
\end{thebibliography}
\end{document}
